% K-means. Subsection of "Clustering".

\subsection{K-means}
\label{sec:clust_kmeans}

Unlike agglomerative clustering, K-means directly partitions the trees into a fixed number $K$ of
clusters. Each cluster $c$ is represented by its centroid $\boldsymbol{\mu}_c$, and the
partition is chosen to minimize the \emph{inertia}, the total squared distance of the trees to
the centroid of their cluster:
\begin{equation}
  J = \sum_{i=1}^{N} \big\|\mathbf{x}_i - \boldsymbol{\mu}_{c(i)}\big\|_2^2 ,
  \label{eq:kmeans}
\end{equation}
where $c(i)$ is the cluster of tree $i$. Lloyd's algorithm \cite{lloyd1982} alternates two
steps: every tree is assigned to its nearest centroid, then every centroid is moved to the mean
of the trees assigned to it. Neither step can increase $J$, so the algorithm converges, but only
to a local minimum that depends on the initial centroids. We therefore initialize the centroids
with k-means++ \cite{arthur2007kmeanspp}, which spreads them out by picking distant trees with
higher probability, and keep the best of 20 runs with different initializations (scikit-learn
implementation \cite{pedregosa2011sklearn}). Because each tree joins its nearest centroid, the
borders between clusters are flat hyperplanes: K-means works best for compact, roughly
spherical clusters.

\paragraph{Choosing the number of clusters}
K-means requires the number of clusters $K$ to be specified in advance. We therefore evaluated
$K=2,\dots,10$ for each airway-tree representation. For each candidate $K$, K-means was fitted
to the 256-dimensional latent codes using k-means++ initialization with 20 restarts, retaining
the solution with the lowest inertia. Two criteria were then evaluated: the inertia, which
measures the within-cluster sum of squared distances, and the silhouette coefficient
\cite{rousseeuw1987silhouettes}, which measures how well each tree is separated from the other
clusters. The mean silhouette was used to assess the separation of the resulting partition.
\begin{equation}
  s(i) = \frac{b(i) - a(i)}{\max\{a(i),\, b(i)\}} \in [-1, 1].
  \label{eq:silhouette}
\end{equation}
The analysis identified $K=3$ as the preferred number of clusters for both the registered and
unregistered airway-tree representations. The final K-means analysis was therefore performed
with three clusters for both representations.